%%%%%%%%%%%%%%%%%%%%%%%%%%%%%%%%%%%%%%%%%%%%%%%%%%%%%%%%%%%%%%%%%%%%%%
% Lớp 8 | Chương IV. Định lí Thalès | Định lí Thalès trong tam giác
% Dạng 8H4?15-1: Tỉ số của hai đoạn thẳng, đoạn thẳng tỉ lệ
%%%%%%%%%%%%%%%%%%%%%%%%%%%%%%%%%%%%%%%%%%%%%%%%%%%%%%%%%%%%%%%%%%%%%%


% nguồn: dan_tu_AI (nhập 16/11/2026)
\begin{bt}%[8H4H15-1]%[ma:d5078332]
	a) Tính tỉ số của hai đoạn thẳng $AB=6\text{ cm}$ và $CD=15\text{ cm}$.\\
	b) Tính tỉ số của hai đoạn thẳng $EF=40\text{ mm}$ và $GH=1{,}2\text{ dm}$.
	\loigiai{
		a) $\dfrac{AB}{CD}=\dfrac{6}{15}=\dfrac25$.\\
		b) Đổi về cùng đơn vị: $EF=40\text{ mm}=4\text{ cm}$, $GH=1{,}2\text{ dm}=12\text{ cm}$.\\
		$\dfrac{EF}{GH}=\dfrac{4}{12}=\dfrac13$.
	}
\end{bt}

% nguồn: dan_tu_AI (nhập 16/11/2026)
\begin{bt}%[8H4H15-1]%[ma:50e55776]
	a) Biết $\dfrac{AB}{CD}=\dfrac34$ và $CD=12\text{ cm}$. Tính độ dài $AB$.\\
	b) Hai đoạn thẳng $AB=3\text{ cm}$, $CD=5\text{ cm}$ có tỉ lệ với hai đoạn thẳng $A'B'=6\text{ cm}$, $C'D'=10\text{ cm}$ không? Vì sao?
	\loigiai{
		a) $AB=\dfrac34\cdot CD=\dfrac34\cdot 12=9\text{ (cm)}$.\\
		b) $\dfrac{AB}{CD}=\dfrac35$ và $\dfrac{A'B'}{C'D'}=\dfrac{6}{10}=\dfrac35$.\\
		Vì $\dfrac{AB}{CD}=\dfrac{A'B'}{C'D'}$ nên hai đoạn thẳng $AB$, $CD$ tỉ lệ với hai đoạn thẳng $A'B'$, $C'D'$.
	}
\end{bt}

% nguồn: dan_tu_AI (nhập 16/11/2026)
\begin{bt}%[8H4H15-1]%[ma:d4549044]
	Cho đoạn thẳng $AB=10\text{ cm}$ và điểm $M$ nằm giữa $A$ và $B$ sao cho $\dfrac{MA}{MB}=\dfrac23$. Tính độ dài các đoạn thẳng $MA$ và $MB$.
	\loigiai{
		Từ $\dfrac{MA}{MB}=\dfrac23$ suy ra $\dfrac{MA}{2}=\dfrac{MB}{3}=\dfrac{MA+MB}{2+3}=\dfrac{AB}{5}=\dfrac{10}{5}=2$.\\
		Vậy $MA=2\cdot 2=4\text{ (cm)}$, $MB=2\cdot 3=6\text{ (cm)}$.
	}
\end{bt}

% nguồn: dan_tu_AI (nhập 16/11/2026)
\begin{bt}%[8H4H15-1]%[ma:2190d17d]
	a) Cho $AB=4\text{ cm}$, $CD=6\text{ cm}$, $EF=10\text{ cm}$, $GH=15\text{ cm}$. Hai đoạn thẳng $AB$, $CD$ có tỉ lệ với hai đoạn thẳng $EF$, $GH$ không?\\
	b) Cho $MN=3\text{ cm}$, $PQ=5\text{ cm}$, $RS=6\text{ cm}$, $UV=9\text{ cm}$. Hai đoạn thẳng $MN$, $PQ$ có tỉ lệ với hai đoạn thẳng $RS$, $UV$ không?
	\loigiai{
		a) $\dfrac{AB}{CD}=\dfrac46=\dfrac23$; $\dfrac{EF}{GH}=\dfrac{10}{15}=\dfrac23$.\\
		Vì $\dfrac{AB}{CD}=\dfrac{EF}{GH}$ nên $AB$, $CD$ tỉ lệ với $EF$, $GH$.\\
		b) $\dfrac{MN}{PQ}=\dfrac35$; $\dfrac{RS}{UV}=\dfrac69=\dfrac23$.\\
		Vì $\dfrac35\ne\dfrac23$ nên $MN$, $PQ$ không tỉ lệ với $RS$, $UV$.
	}
\end{bt}

% nguồn: dan_tu_AI (nhập 16/11/2026)
\begin{bt}%[8H4H15-1]%[ma:e9bbc550]
	Trên một bản đồ có tỉ lệ $1:50\,000$, khoảng cách giữa hai điểm $A$ và $B$ đo được là $6\text{ cm}$. Tính khoảng cách thực tế giữa hai địa điểm đó (đơn vị: km).
	\loigiai{
		Tỉ lệ bản đồ $1:50\,000$ nghĩa là tỉ số giữa độ dài trên bản đồ và độ dài thực tế bằng $\dfrac{1}{50\,000}$.\\
		Khoảng cách thực tế: $6\cdot 50\,000=300\,000\text{ (cm)}=3\text{ (km)}$.
	}
\end{bt}

% nguồn: dan_tu_AI (nhập 16/11/2026)
\begin{bt}%[8H4V15-1]%[ma:a4998470]
	Cho đoạn thẳng $AB=20\text{ cm}$.\\
	a) Tìm điểm $M$ nằm giữa $A$ và $B$ sao cho $\dfrac{MA}{MB}=\dfrac32$. Tính $MA$, $MB$.\\
	b) Tìm điểm $N$ nằm trên tia đối của tia $BA$ sao cho $\dfrac{NA}{NB}=\dfrac32$. Tính $NB$.\\
	c) Tính độ dài đoạn thẳng $MN$.
	\loigiai{
		a) $\dfrac{MA}{3}=\dfrac{MB}{2}=\dfrac{MA+MB}{5}=\dfrac{20}{5}=4$, suy ra $MA=12\text{ cm}$, $MB=8\text{ cm}$.\\
		b) $N$ nằm trên tia đối của tia $BA$ nên $B$ nằm giữa $A$ và $N$, do đó $NA=NB+20$.\\
		$\dfrac{NB+20}{NB}=\dfrac32\Leftrightarrow 2(NB+20)=3NB\Leftrightarrow NB=40\text{ (cm)}$.\\
		c) $B$ nằm giữa $M$ và $N$ nên $MN=MB+BN=8+40=48\text{ (cm)}$.
	}
\end{bt}

% nguồn: dan_tu_AI (nhập 18/11/2026)
\begin{ex}%[8H4N15-1]%[ma:195dd86b]
	Cho hai đoạn thẳng $AB = 15\text{ cm}$ và $CD = 20\text{ cm}$. Tỉ số của hai đoạn thẳng $AB$ và $CD$ là
	\choice
	{$\dfrac{4}{3}$}
	{\True $\dfrac{3}{4}$}
	{$\dfrac{3}{7}$}
	{$\dfrac{7}{3}$}
	\loigiai{
		Tỉ số của hai đoạn thẳng là $\dfrac{AB}{CD} = \dfrac{15}{20} = \dfrac{3}{4}$.
	}
\end{ex}

% nguồn: dan_tu_AI (nhập 18/11/2026)
\begin{ex}%[8H4N15-1]%[ma:cf51d949]
	Đoạn thẳng $MN$ tỉ lệ với đoạn thẳng $PQ$ theo tỉ số $k=2$. Biết $PQ = 4\text{ cm}$, độ dài đoạn thẳng $MN$ là
	\choice
	{$2\text{ cm}$}
	{$6\text{ cm}$}
	{\True $8\text{ cm}$}
	{$16\text{ cm}$}
	\loigiai{
		Ta có $\dfrac{MN}{PQ} = 2 \Rightarrow MN = 2 \cdot PQ = 2 \cdot 4 = 8\text{ cm}$.
	}
\end{ex}

% nguồn: dan_tu_AI (nhập 18/11/2026)
\begin{ex}%[8H4H15-1]%[ma:17d29d97]
	Tỉ số của đoạn thẳng $AB = 45\text{ cm}$ và đoạn thẳng $CD = 1{,}5\text{ m}$ là
	\choice
	{$30$}
	{\True $\dfrac{3}{10}$}
	{$\dfrac{10}{3}$}
	{$3$}
	\loigiai{
		Đổi $1{,}5\text{ m} = 150\text{ cm}$. Tỉ số là $\dfrac{AB}{CD} = \dfrac{45}{150} = \dfrac{3}{10}$.
	}
\end{ex}

% nguồn: dan_tu_AI (nhập 18/11/2026)
\begin{ex}%[8H4N15-1]%[ma:faca785e]
	Hai đoạn thẳng $AB = 6\text{ cm}$ và $CD = 8\text{ cm}$. Các khẳng định sau đúng hay sai?
	\choiceTF
	{Tỉ số $\dfrac{AB}{CD} = \dfrac{4}{3}$}
	{\True Tỉ số $\dfrac{CD}{AB} = \dfrac{4}{3}$}
	{Tỉ số của hai đoạn thẳng phụ thuộc vào đơn vị đo}
	{\True Nếu $MN = 3\text{ cm}, PQ = 4\text{ cm}$ thì $AB, CD$ tỉ lệ với $MN, PQ$}
	\loigiai{
		\begin{itemchoice}
			\itemch Sai. $\dfrac{AB}{CD} = \dfrac{6}{8} = \dfrac{3}{4}$.
			\itemch Đúng. $\dfrac{CD}{AB} = \dfrac{8}{6} = \dfrac{4}{3}$.
			\itemch Sai. Tỉ số của hai đoạn thẳng không phụ thuộc vào đơn vị đo (phải cùng đơn vị).
			\itemch Đúng. $\dfrac{AB}{CD} = \dfrac{3}{4}$ và $\dfrac{MN}{PQ} = \dfrac{3}{4}$ nên chúng tỉ lệ.
		\end{itemchoice}
	}
\end{ex}

% nguồn: dan_tu_AI (nhập 18/11/2026)
\begin{bt}%[8H4N15-1]%[ma:43669a64]
	Hai đoạn thẳng $MN = 15\text{ cm}$ và $PQ = 25\text{ cm}$. Tỉ số $\dfrac{MN}{PQ}$ là một phân số tối giản $\dfrac{a}{b}$. Tính $a+b$.
	\shortans{8}
	\loigiai{
		$\dfrac{MN}{PQ} = \dfrac{15}{25} = \dfrac{3}{5}$. Vậy $a=3, b=5 \Rightarrow a+b = 8$.
	}
\end{bt}

% nguồn: dan_tu_AI (nhập 18/11/2026)
\begin{bt}%[8H4N15-1]%[ma:6e6bb7ee]
	Đoạn $AB=10\text{ cm}$ và đoạn $CD=2\text{ dm}$. Tính tỉ số $\dfrac{AB}{CD}$.
	\shortans{0,5}
	\loigiai{
		Đổi $2\text{ dm} = 20\text{ cm}$. $\dfrac{AB}{CD} = \dfrac{10}{20} = 0{,}5$.
	}
\end{bt}

% nguồn: dan_tu_AI (nhập 18/11/2026)
\begin{bt}%[8H4N15-1]%[ma:a8c7865e]
	Cho 4 đoạn thẳng $a = 4\text{ cm}, b = 6\text{ cm}, c = 10\text{ cm}, d = 15\text{ cm}$. Hỏi $a, b$ có tỉ lệ với $c, d$ không? Vì sao?
	\loigiai{
		Ta có:
		\[ \dfrac{a}{b} = \dfrac{4}{6} = \dfrac{2}{3} \]
		\[ \dfrac{c}{d} = \dfrac{10}{15} = \dfrac{2}{3} \]
		Vì $\dfrac{a}{b} = \dfrac{c}{d}$ nên hai đoạn thẳng $a, b$ tỉ lệ với $c, d$.
	}
\end{bt}

% nguồn: 2026-11-18_104705_tuong-tu-theo-mau (nhập 18/11/2026)
\begin{bt}%[8H4N15-1]%[ma:1a489e49]
% nguon: AI soạn tương tự câu mẫu a8c7865e
% mau: a8c7865e
Cho 4 đoạn thẳng $a = 8\text{ cm}, b = 12\text{ cm}, c = 14\text{ cm}, d = 21\text{ cm}$. Hỏi $a, b$ có tỉ lệ với $c, d$ không? Vì sao?
	\loigiai{
		Ta có:\\
		\[ \dfrac{a}{b} = \dfrac{8}{12} = \dfrac{2}{3} \]\\
		\[ \dfrac{c}{d} = \dfrac{14}{21} = \dfrac{2}{3} \]\\
		Vì $\dfrac{a}{b} = \dfrac{c}{d}$ nên hai đoạn thẳng $a, b$ tỉ lệ với $c, d$.
	}
\end{bt}

% nguồn: 2026-11-18_104705_tuong-tu-theo-mau (nhập 18/11/2026)
\begin{bt}%[8H4N15-1]%[ma:5a24f6b7]
% nguon: AI soạn tương tự câu mẫu a8c7865e
% mau: a8c7865e
Cho 4 đoạn thẳng $a = 9\text{ cm}, b = 15\text{ cm}, c = 12\text{ cm}, d = 20\text{ cm}$. Hỏi $a, b$ có tỉ lệ với $c, d$ không? Vì sao?
	\loigiai{
		Ta có:\\
		\[ \dfrac{a}{b} = \dfrac{9}{15} = \dfrac{3}{5} \]\\
		\[ \dfrac{c}{d} = \dfrac{12}{20} = \dfrac{3}{5} \]\\
		Vì $\dfrac{a}{b} = \dfrac{c}{d}$ nên hai đoạn thẳng $a, b$ tỉ lệ với $c, d$.
	}
\end{bt}

% nguồn: 2026-11-18_104705_tuong-tu-theo-mau (nhập 18/11/2026)
\begin{bt}%[8H4H15-1]%[ma:18bae821]
% nguon: AI soạn tương tự câu mẫu a8c7865e
% mau: a8c7865e nang_cao
	Cho ba đoạn thẳng $a = 6\text{ cm}$, $b = 9\text{ cm}$ và đoạn thẳng $d = 12\text{ cm}$. Tìm độ dài đoạn thẳng $c$ để hai đoạn thẳng $a, b$ tỉ lệ với $c, d$.
	\loigiai{
		Để hai đoạn thẳng $a, b$ tỉ lệ với $c, d$ thì ta phải có $\dfrac{a}{b} = \dfrac{c}{d}$.\\
		Ta có $\dfrac{a}{b} = \dfrac{6}{9} = \dfrac{2}{3}$.\\
		Suy ra $\dfrac{c}{d} = \dfrac{2}{3}$, tức là $\dfrac{c}{12} = \dfrac{2}{3}$.\\
		Do đó $c = \dfrac{12 \cdot 2}{3} = 8$ (cm).\\
		Thử lại: $\dfrac{a}{b} = \dfrac{6}{9} = \dfrac{2}{3}$ và $\dfrac{c}{d} = \dfrac{8}{12} = \dfrac{2}{3}$, đúng.\\
		Vậy $c = 8\text{ cm}$.
	}
\end{bt}

% nguồn: De_AI_Toan8_cuoi_hoc_ki_i_1 (nhập 21/11/2026)
\begin{bt}%[8H4H15-1]%[ma:78ba965e]
	Cho các đoạn thẳng $AB = 6\text{ cm}$, $CD = 8\text{ cm}$, $EF = 9\text{ cm}$, $GH = 12\text{ cm}$.
a) Tính các tỉ số $\dfrac{AB}{CD}$ và $\dfrac{EF}{GH}$.
b) Hai đoạn thẳng $AB$, $CD$ có tỉ lệ với hai đoạn thẳng $EF$, $GH$ không? Vì sao?
c) Tìm độ dài đoạn thẳng $MN$ biết $\dfrac{MN}{AB} = \dfrac{5}{3}$.
	\loigiai{
		a) Ta có $\dfrac{AB}{CD} = \dfrac{6}{8} = \dfrac{3}{4}$ và $\dfrac{EF}{GH} = \dfrac{9}{12} = \dfrac{3}{4}$. \\
b) Vì $\dfrac{AB}{CD} = \dfrac{EF}{GH} = \dfrac{3}{4}$ nên hai đoạn thẳng $AB$, $CD$ tỉ lệ với hai đoạn thẳng $EF$, $GH$. \\
c) Từ $\dfrac{MN}{AB} = \dfrac{5}{3}$, suy ra $MN = \dfrac{5}{3} \cdot AB = \dfrac{5}{3} \cdot 6 = 10\text{ (cm)}$.
	}
\end{bt}

% nguồn: De_AI_Toan8_cuoi_hoc_ki_i_1 (nhập 21/11/2026)
\begin{bt}%[8H4H15-1]%[ma:1d3d1fb5]
	Cho đoạn thẳng $AB = 21\text{ cm}$. Điểm $M$ nằm trên đoạn thẳng $AB$ sao cho $\dfrac{MA}{MB} = \dfrac{3}{4}$.
a) Tính độ dài các đoạn thẳng $MA$ và $MB$.
b) Lấy điểm $N$ nằm trên tia đối của tia $AB$ sao cho $\dfrac{NA}{NB} = \dfrac{3}{4}$. Tính độ dài các đoạn thẳng $NA$ và $NB$.
	\loigiai{
		a) Vì $M$ nằm trên đoạn thẳng $AB$ nên $MA + MB = AB = 21\text{ (cm)}$. \\
Từ $\dfrac{MA}{MB} = \dfrac{3}{4}$, áp dụng tính chất dãy tỉ số bằng nhau ta có: \\
$\dfrac{MA}{3} = \dfrac{MB}{4} = \dfrac{MA + MB}{3 + 4} = \dfrac{21}{7} = 3$. \\
Suy ra $MA = 3 \cdot 3 = 9\text{ (cm)}$ và $MB = 4 \cdot 3 = 12\text{ (cm)}$. \\
b) Vì $N$ nằm trên tia đối của tia $AB$ nên điểm $A$ nằm giữa $N$ và $B$, do đó $NB - NA = AB = 21\text{ (cm)}$. \\
Từ $\dfrac{NA}{NB} = \dfrac{3}{4}$, suy ra $\dfrac{NA}{3} = \dfrac{NB}{4} = \dfrac{NB - NA}{4 - 3} = \dfrac{21}{1} = 21$. \\
Vậy $NA = 3 \cdot 21 = 63\text{ (cm)}$ và $NB = 4 \cdot 21 = 84\text{ (cm)}$.
	}
\end{bt}
